\section{Simulation: what holds, and where it stops}\label{sec:simulation}

Three questions, and the third is the one that decides whether the procedure is
worth using.

\subsection{Against a comparator sharing the target, information and level}

The direct comparator is a normal-theory prediction interval for the \emph{same}
new population's latent value, built from the \emph{same} structure confidence
set and the \emph{same} certified scale bound, spending the identical budget
$\alpha_0+\eta_\gamma+\eta_t=0.10$: $A_U=\bar S/\chi^2_{K,\alpha_2}$, the exact
chi-square limit rather than a Satterthwaite or \citet{graybill1980}
approximation \citep{satterthwaite1946}, and
$\pm z_{1-\alpha_1/2}\sqrt{A_U}$ with $\alpha_1+\alpha_2=\alpha_0$. The two arms
differ in exactly one place.

\begin{center}
\begin{tabular}{lcccc}
\toprule
$K$ & proposed cov. & comparator cov. & \textbf{proposed narrower by} & split\\
\midrule
60 & 0.957--0.978 & 0.982--0.996 & 22.2\% & scanned in its own favour\\
60 & --- & --- & \textbf{25.9\%} & fixed in advance\\
200 & 0.957--0.970 & 0.983--0.988 & 13.9\% & scanned in its own favour\\
200 & --- & --- & \textbf{20.4\%} & fixed in advance\\
\bottomrule
\end{tabular}
\end{center}

All against a guaranteed 0.90. The pre-registered split is the primary
comparison; the scanned one gives the comparator its best budget allocation and
is the smaller number, so \textbf{we quote 13.9 to 22.2\% throughout}. Two
readings follow. Both arms over-cover, but the proposal is \emph{both} narrower
and closer to the nominal level, at 0.957--0.978 against 0.982--0.996: the width
difference is not bought by spending coverage. And the comparator is one valid
construction, not a class and not the best possible---the registered prediction
before this experiment was that it would win, and it did not.

\paragraph{A disagreement between two of our own runs.} On a tight scale limit
the normal-theory band is narrower by 2 to 22\% in eight cells; on a loose pivot
limit an earlier run found it wider by 7 to 14\%. We recorded the second as an
error and superseded it. Consequence~2 of Section~\ref{sec:sensitivity} says the
two are the two sides of $\lambda^\star$ and neither need be wrong. We report
this as a reconciliation and not as evidence: the runs differ in more than the
tightness of the limit, and an experiment that varies only $\lambda$ and locates
the crossing is one we have not done.

\subsection{The regimes where it stops}

Under finite-population stratified cluster sampling with a pre-specified
sample-size structure variable, the picture divides not by design share but by
what else changes with it.

\begin{center}
\begin{tabular}{lcc}
\toprule
& degrees of freedom 12--28 & degrees of freedom 4--8\\
\midrule
structure $R^2$ & 0.59--0.85 & 0.05--0.71\\
zero variance estimates & none & up to 17\%\\
\textbf{containment} & \textbf{0.955--0.992} & \textbf{0.681--0.987}\\
latent coverage & 0.958--0.971 & 0.944--0.975\\
band vs.\ structure-free & 0.883--0.937 & 0.836--0.928\\
\bottomrule
\end{tabular}
\end{center}

At matched $\bar D/A\in\{0.55,0.75\}$ the two columns differ in degrees of
freedom, in the rate of exactly-zero variance estimates and in the fit of the
structure model---\textbf{the same noise ratio is not the same inference
problem}. In the first column the procedure holds its containment requirement and
is 6.3 to 11.7\% narrower than a structure-free construction, capturing 0.56--0.74
of the available narrowing against 0.23--0.30. In the second it loses containment
at a tail quantile and its width figures are not usable.

Supplying the \emph{true} structure raises containment there from 0.68 to 0.82
and from 0.77 to 0.89 and does not reach 0.95. Two conclusions and only two:
approximating the variance structure affects the outcome, and improving that step
alone would not repair the procedure. We do not apportion the residue among
normality, the $\chi^2$ law and independence; this design breaks all three at
once and the experiment does not separate them.

\paragraph{We report the verified scenario range, not a usage threshold.} Writing
``degrees of freedom above 12 and structure $R^2$ above 0.6'' as an operating rule
would be wrong three times: those numbers co-occurred in these scenarios rather
than being derived as thresholds; the structure $R^2$ is computed from the
\emph{true} design variances and is unavailable in real data; and nothing here
proves normality or the $\chi^2$ law under any complex design.

\subsection{Honest accounting of the gain}

At the regional design share the \emph{valid} interval is 0.9 to 6.6\% narrower
than not correcting at all, capturing 6 to 31\% of what the oracle achieves. The
reason is the guarantee itself: the argument needs $A$ bounded above and each
$D_c$ bounded below, and both drive the correction toward zero. Measuring where
the width goes shows the simultaneity correction is only 6 to 22\% of that loss,
so replacing it alone would not repair matters---which is why the construction
bounds the ratio rather than the components.
